\documentclass[11pt,a4paper]{article}
\usepackage[T1]{fontenc}
\usepackage[utf8]{inputenc}
\usepackage[margin=2.4cm]{geometry}
\usepackage{enumitem}
\setlist{itemsep=2pt,topsep=3pt}
\usepackage{booktabs}
\usepackage{array}
\usepackage{xcolor}
\definecolor{accent}{HTML}{3F5B74}
\usepackage[colorlinks=true,linkcolor=accent,urlcolor=accent]{hyperref}
\usepackage{titlesec}
\titleformat{\section}{\large\bfseries\color{accent}}{\thesection.}{0.5em}{}
\titleformat{\subsection}{\bfseries}{\thesubsection}{0.5em}{}
\setlength{\parindent}{0pt}
\setlength{\parskip}{0.5em}

\title{\vspace{-1.5cm}\textbf{Lecture Summariser --- Proposal}}
\author{}
\date{Updated 2026-06-18}

\begin{document}
\maketitle
\vspace{-1.2em}

\noindent\textbf{One line.} A fully-local, offline tool that turns a lecture \emph{slide PDF} into
structured, textbook-grounded \LaTeX/PDF revision notes --- no external APIs, nothing leaves the machine.

\begin{quote}\small
This proposal was rewritten on 2026-06-18 to describe the product that actually ships. The earlier
roadmap (converted from \texttt{LLM lecture note taker.docx}, 2026-06-02) framed multimodal image/diagram
grounding as the ``core novelty''; that component was never built, while the working slides-to-notes
generator was undersold by its own spec. This version makes the shipped tool the headline and demotes the
unbuilt parts to an explicit, optional future track. The previous version is archived under
\texttt{\_archive/}.
\end{quote}

\section{What it is}

The tool ingests a course's lecture slide decks and produces clean, styled revision notes a student can
revise from. It is deliberately \textbf{fully local and offline}: a small open-weight LLM via
\href{https://ollama.com}{Ollama} (default \texttt{qwen2.5}, swappable) plus a local sentence-embedding
retriever over a course textbook. No paid or cloud APIs are used at inference, and no slide, transcript, or
note ever leaves the user's machine. This offline, privacy-preserving, \LaTeX-native niche is the project's
genuine differentiation: mainstream study tools (e.g. NotebookLM, Otter, Plaud) are cloud-hosted and do not
emit compilable \LaTeX.

\textbf{Primary input:} a slide deck (PDF). \textbf{Optional input:} a plain-text transcript/caption file,
used as extra retrieval context (not a co-equal source). \textbf{Output:} a compilable \texttt{.tex} file
and its rendered PDF, with definitions, key facts, comparison tables, worked examples and traps in a muted,
consistent style.

\section{Architecture (the live pipeline)}

For each deck, \texttt{run\_extraction.run\_pipeline}:

\begin{enumerate}[leftmargin=1.4em]
\item \textbf{Extract \& clean.} Pull slide text with PyMuPDF; strip placeholder glyphs that break \LaTeX.
\item \textbf{Plan an outline.} Group slides by their \texttt{Section: Subtitle} header into one topic each
  (merging split themes, dropping syllabus/objectives/feedback slides). Every topic becomes its own
  \texttt{\textbackslash section}, so none can be silently dropped.
\item \textbf{Per topic, a staged Content $\rightarrow$ Review $\rightarrow$ Structure pipeline}, each stage
  a small focused local-LLM call:
  \begin{itemize}
  \item \textbf{Scope-linked grounding:} per-slide retrieval from the local textbook index, then a local-LLM
    \emph{scope gate} that keeps only passages whose primary subject is the topic. The ``different subject''
    contrast is derived from the deck's own other topics, so the gate is not hard-wired to one course.
  \item \textbf{Protected spans:} the slide's own worked examples / vectors / similarity matrices / tables are
    extracted verbatim, frozen as an opaque \texttt{[[EXAMPLES]]} token, deduped across topics, and
    substituted back as real \LaTeX{} only at the very end --- so no stage can re-author or break them.
  \item \textbf{Content:} grounded numbered points (coverage + correctness).
  \item \textbf{Review:} the model fact-checks each point against the source.
  \item \textbf{Independent NLI check:} a cross-encoder entailment model (a different architecture and
    failure-mode than the writer) scores each point against the source and \emph{drops points the source
    contradicts} --- the swapped-definition / invented-number class of error that same-model self-review
    misses.
  \item \textbf{Structure (deterministic):} the model emits a \emph{structured JSON} description of the
    section (typed blocks --- definitions, comparison tables, fact cards, traps, prose); a Python renderer
    turns that into guaranteed-clean \LaTeX. Comparisons always become real tables, every box is balanced
    and correctly labelled, nothing nests. Formatting is taken off the small model, so \textbf{structure can
    no longer be the limiting factor}.
  \end{itemize}
\item \textbf{Deterministic dedup \& cleanup}; then render via a styled preamble and \texttt{xelatex}.
\end{enumerate}

\section{Status and results}

Working end-to-end. Notes are scored 1--10 each on Coverage / Correctness / Structure by an independent
frontier-model judge reading the source slides and the generated notes side by side (reproducible harness in
\texttt{eval/}). A blind 3-judge A/B on the held-out L3 deck (out of 30) shows both the deterministic
structure renderer + NLI check and a newer model improving output: OLD (LLM-authored LaTeX, no NLI)
\textbf{14.6} $\rightarrow$ NEW (deterministic + NLI, qwen2.5-7b) \textbf{15.6} $\rightarrow$
NEW + qwen3-8b \textbf{16.0}. Correctness rose $3.0\to4.3\to5.0$ and Structure $4.3\to5.0$; objective metrics
show comparison tables $3\to6$ on L3 with zero LaTeX defects. Held-out L5 generalises (15.0). The cost is
about one point of Coverage (the NLI filter and renderer trade breadth for fidelity).

\textbf{Honest limitations.} (i) \emph{Factual correctness} is still the ceiling: the NLI check catches the
\emph{contradiction} class but keeps neutral content, so fabrication-by-omission (an invented section from a
near-empty slide) survives --- the next lever is a groundedness/\emph{support} gate, plus a newer/bigger
model. (ii) \emph{Evaluation} is stronger (a 3-judge panel + objective metrics) but still NLP-course decks
only --- no non-NLP deck exists yet to test true generalisation. (iii) \emph{Course coverage:} only one
textbook is indexed (Jurafsky \& Martin); a different course needs its own book.

\section{Roadmap}

\emph{Done (2026-06-18):} deterministic structure renderer, independent NLI fact-check, de-hardcoded scope
gate, the newer-model A/B (qwen3-8b), and a multi-judge eval harness --- see Status and results.

\subsection*{Active (next, the shipped tool)}
\begin{itemize}
\item \textbf{Groundedness / support gate} --- the now-dominant correctness hole: drop a block when nothing
  in the source \emph{entails} it (not only when it is contradicted), and emit no section for a slide range
  with no extractable text. Upgrade the NLI checker toward a dedicated model such as MiniCheck.
\item \textbf{Consolidate comparisons} --- the renderer guarantees a valid table, but the model must emit a
  single comparison block; tighten the structure prompt so multi-way comparisons stop fragmenting.
\item \textbf{Recover the $\sim$1-point coverage cost} (tune the NLI thresholds) and run one non-NLP deck once
  one is supplied.
\item A stronger local retriever than \texttt{all-MiniLM-L6-v2} (e.g. EmbeddingGemma / BGE / GTE).
\end{itemize}

\subsection*{Optional / future (only if it earns its keep)}
\begin{itemize}
\item \textbf{Figure grounding.} Local, on-laptop multimodal models (e.g. Qwen2.5-VL-7B) make extracting and
  explaining slide figures feasible. The cheapest viable slice is placing already-extracted slide images
  inline under their topic with no model narration; full per-figure captioning is deferred because unreliable
  narration is worse than none.
\item Borrow source-citation (NotebookLM-style page references) so each claim is traceable to a slide.
\end{itemize}

\subsection*{Cut (explicitly out of scope)}
Flashcard generation, multi-lecture stitching, and direct \texttt{.pptx} input are dropped as scope creep
for a single-user tool. An internet/web-research enrichment agent is dropped outright --- it would violate
the fully-local/offline premise that is the project's defining constraint.

\section{Core technical challenges}
\begin{enumerate}[leftmargin=1.4em]
\item Keeping notes faithful to the slides on a small local model (the correctness ceiling).
\item Grounding without over-retrieval: keeping only textbook passages genuinely on-topic.
\item Clean, compilable \LaTeX{} across varied slide layouts --- now owned by the deterministic renderer.
\item Generalising beyond the one course the tool was tuned on.
\end{enumerate}

\end{document}
